\section{Exploratory results \texorpdfstring{\screen}{[screen]}}\label{sec:explore}

Everything in this section was not preregistered, or uses all \NScenarios\ scenarios (training and test
together). It explains and qualifies the confirmatory results; it does not confirm anything on its own.

\subsection{Where the cost comes from}\label{sec:composition}

\begin{figure}[htbp]
\centering
\pgfplotslegendfromname{complegend}\par\smallskip
% Where the money goes: mean tools-normalized cost per session, split by token class.
\begin{tikzpicture}
\begin{axis}[benchmark, xbar stacked, bar width=10pt, scale only axis, width=10.5cm, height=6.4cm,
  xmin=0, xmax=\CompMax, ymin=-0.6, ymax=6.6, ytick={0,...,6},
  yticklabels from table={generated/data/composition.dat}{label},
  xlabel={Mean cost per session (US dollars)}, xmajorgrids,
  legend to name=complegend, legend columns=3, legend cell align=left,
  legend style={/tikz/every even column/.append style={column sep=10pt}}]
\addplot[fill=okgray!50, draw=black!60] table[x=input, y=idx]{generated/data/composition.dat};
\addlegendentry{uncached input}
\addplot[fill=okblue!45, draw=black!60] table[x=read, y=idx]{generated/data/composition.dat};
\addlegendentry{cache read}
\addplot[fill=okorange!80, draw=black!60, postaction={pattern=north east lines, pattern color=white}]
  table[x=write, y=idx]{generated/data/composition.dat};
\addlegendentry{cache write (new content)}
\addplot[fill=black!75, draw=black!60, postaction={pattern=crosshatch dots, pattern color=white}]
  table[x=rebuild, y=idx]{generated/data/composition.dat};
\addlegendentry{rebuild after a switch}
\addplot[fill=okgreen!70, draw=black!60, postaction={pattern=dots, pattern color=white}]
  table[x=output, y=idx]{generated/data/composition.dat};
\addlegendentry{output}
\end{axis}
\end{tikzpicture}

\caption{Mean cost per session, split by what was paid for (all splits, cost-valid sessions,
tools-normalized). ``Rebuild after a switch'' is the extra paid for re-writing history to a new model's
cache.}
\label{fig:composition}
\howtoread{Each bar is one kind of session; its length is the average bill and its segments show what the
bill was for. Compare the three Fable bars with each other, and the three Opus bars with each other. On
Fable, routing shrinks the large hatched ``cache write'' segment. On Opus, routing lengthens the blue
``cache read'' segment.}
\end{figure}

\begin{table}[htbp]
\centering\small
\caption{The numbers behind \cref{fig:composition}: dollars per session by token class, main-loop requests per
session, and cache-read tokens per session (millions).}
\label{tab:composition}
\begin{tabular*}{\textwidth}{@{\extracolsep{\fill}}l r r r r r r r r@{}}
\toprule
 & & & \multicolumn{4}{c}{of which (\$ per session)} & & \\
\cmidrule(lr){4-7}
Session kind & sessions & \$/session & cache read & cache write & rebuild & output & requests & tokens read \\
\midrule
Fable: plain host & 140 & 5.31 & 1.06 & 2.82 & 0.00 & 1.27 & 46.0 & 4.22\,M \\
Fable: shipped & 140 & 3.56 & 1.37 & 0.97 & 0.61 & 0.49 & 50.6 & 4.68\,M \\
Fable: sticky & 140 & 3.33 & 1.34 & 1.40 & 0.00 & 0.45 & 50.2 & 4.58\,M \\
Sonnet only (control) & 140 & 3.22 & 1.77 & 0.96 & 0.00 & 0.44 & 57.4 & 5.90\,M \\
Opus: plain host & 140 & 2.13 & 0.62 & 1.04 & 0.00 & 0.40 & 36.4 & 3.07\,M \\
Opus: shipped & 140 & 2.71 & 1.28 & 0.77 & 0.29 & 0.32 & 48.9 & 4.43\,M \\
Opus: sticky & 140 & 2.61 & 1.28 & 0.95 & 0.00 & 0.33 & 48.5 & 4.43\,M \\
\bottomrule
\end{tabular*}

\end{table}

\Cref{fig:composition,tab:composition} show why the two hosts disagree.
\begin{itemize}
  \item \textbf{On Fable, routing cuts the most expensive line.} A plain Fable session paid \$\CompFableAnchorWrite\
    of its \$\CompFableAnchorTotal\ for cache writes (\WriteShareFableAnchor\,\%), and a Sonnet cache write costs
    \SonnetFableWritePriceX\X\ a Fable one. Moving the session to Sonnet cuts that line by more than the extra reads
    add.
  \item \textbf{On Opus, routing makes the session longer and its reads dearer.} A sticky session on the Opus
    host made \CompOpusStickyReq\ requests against the plain host's \CompOpusAnchorReq\
    (\ReqRatioOpusSticky\X), read \CompOpusStickyReadTok\ million cached tokens against
    \CompOpusAnchorReadTok\ million, and paid \SonnetOpusReadPriceX\X\ Opus's price for each read: cache
    reads rose from \$\CompOpusAnchorRead\ to \$\CompOpusStickyRead\ per session.
  \item \textbf{Switching adds a rebuild.} The shipped arm re-chooses its model at the start of each turn and
    switched models \SwitchFable\ times per session on average, always between turns, in \SwitchShareFable\,\% of Fable sessions; rebuilds were \RebuildFable\,\% of its cost on
    Fable and \RebuildOpus\,\% on Opus. Sticky never switches (\SwitchStickyFable\ switches per session),
    so it never pays this line.
\end{itemize}

\subsection{By task type, pause pattern and length}\label{sec:slices}

\begin{figure}[htbp]
\centering
\pgfplotslegendfromname{savingslegendfable}\par\smallskip
% Savings per 1,000 sessions by task type and arm (empirical, all splits), 90% bootstrap ranges.
% #1 = host file stem, #2 = title, #3 xmin, #4 xmax
\newcommand{\savingspanel}[4]{%
\begin{tikzpicture}
\begin{axis}[benchmark, xbar, bar width=4.2pt, scale only axis, width=11cm, height=6.2cm,
  xmin=#3, xmax=#4, ymin=-0.6, ymax=6.6, ytick={0,...,6},
  yticklabels from table={generated/data/savings-#1.dat}{label},
  y tick label style={font=\small},
  xlabel={Dollars saved per 1,000 sessions (negative: extra cost)},
  title={#2}, title style={font=\small\bfseries},
  xmajorgrids, enlarge y limits=false,
  legend to name=savingslegend#1, legend columns=3, legend cell align=left,
  legend style={/tikz/every even column/.append style={column sep=10pt}},
  error bars/error bar style={black!70, line width=0.6pt}]
\draw[black!70, line width=0.8pt] (axis cs:0,-0.6) -- (axis cs:0,6.6);
\addplot[fill=okblue, draw=black!60, error bars/.cd, x dir=both, x explicit]
  table[x=vSticky, y=idx, x error minus=mSticky, x error plus=pSticky]{generated/data/savings-#1.dat};
\addlegendentry{sticky (choose once)}
\addplot[fill=okorange!80, draw=black!60, postaction={pattern=north east lines, pattern color=white},
  error bars/.cd, x dir=both, x explicit]
  table[x=vShipped, y=idx, x error minus=mShipped, x error plus=pShipped]{generated/data/savings-#1.dat};
\addlegendentry{shipped (re-chooses each turn)}
\addplot[fill=okgreen!80, draw=black!60, postaction={pattern=dots, pattern color=white},
  error bars/.cd, x dir=both, x explicit]
  table[x=vSonnet, y=idx, x error minus=mSonnet, x error plus=pSonnet]{generated/data/savings-#1.dat};
\addlegendentry{Sonnet only}
\end{axis}
\end{tikzpicture}}

\savingspanel{fable}{Fable host}{\SavMinFable}{\SavMaxFable}\par\medskip
\savingspanel{opus}{Opus host}{\SavMinOpus}{\SavMaxOpus}
\caption{Dollars saved per 1,000 sessions by task type and arm, measured on all \NScenarios\ scenarios, with
90\,\% scenario-bootstrap ranges. The two panels have different scales.}
\label{fig:savings}
\howtoread{Bars to the right of the zero line are savings; bars to the left are extra costs. On Fable every
bar is to the right. On Opus most bars are to the left, except for the short read-heavy task types (docs,
explain), which are small and rest on \TaskExplainN\ and \TaskDocsN\ scenarios: treat them as hints, not findings.
None of the \RvDocsTrain\ docs scenarios is in the test split, so the docs results rest on training data only.}
\end{figure}

\begin{table}[htbp]
\centering\small
\caption{Measured saving per 1,000 sessions (US dollars) by task type, host and arm, all splits. Ranges are in
\cref{tab:per1000full}.}
\label{tab:per1000}
\begin{tabular*}{\textwidth}{@{\extracolsep{\fill}}l r r r r r r r@{}}
\toprule
 & & \multicolumn{3}{c}{Fable host: saving per 1,000 sessions (\$)} & \multicolumn{3}{c}{Opus host} \\
\cmidrule(lr){3-5}\cmidrule(l){6-8}
Task type & scenarios & sticky & shipped & sonnet & sticky & shipped & sonnet \\
\midrule
bugfix & 18 & 1{,}346 & 797 & 871 & $-$1{,}037 & $-$1{,}240 & $-$1{,}948 \\
feature & 13 & 2{,}183 & 2{,}019 & 2{,}895 & $-$574 & $-$490 & $-$1{,}099 \\
mixed & 28 & 2{,}020 & 1{,}941 & 2{,}405 & $-$303 & $-$401 & $-$815 \\
docs & 4 & 2{,}615 & 2{,}263 & 2{,}246 & 264 & 162 & $-$129 \\
explain & 3 & 2{,}315 & 2{,}244 & 2{,}367 & 341 & 235 & 242 \\
review & 4 & 3{,}007 & 3{,}011 & 2{,}389 & $-$425 & $-$642 & $-$1{,}193 \\
\midrule all tasks & 70 & 1{,}980 & 1{,}754 & 2{,}090 & $-$489 & $-$588 & $-$1{,}096 \\
\bottomrule
\end{tabular*}

\end{table}

Across all tasks, per 1,000 sessions, sticky saved \$\SavFableSticky\ on Fable (90\,\% range
\$\SavFableStickyLo--\$\SavFableStickyHi), shipped \$\SavFableShipped\ and Sonnet \$\SavFableSonnet. On Opus
the same arms \emph{lost} \$\SavOpusSticky, \$\SavOpusShipped\ and \$\SavOpusSonnet. These dollar figures
assume this campaign's mix of tasks and session lengths. They are mean-difference (spend) measures, not
geometric means: for Fable sticky, \$\SavFableSticky\ per 1,000 is exactly 1,000 times the drop in mean cost from
\$\RvSpendAnchor\ to \$\RvSpendSticky. The largest Fable sticky saving by task type is \RvBestTaskFableSticky\
(\$\RvSavReviewFableSticky), ahead of docs (\$\RvSavDocsFableSticky). \Cref{tab:tasksplit} shows how few
scenarios some task types have, and that docs has none in the test split. \Cref{tab:jointtask} sets these savings
next to the quality differences.

\begin{table}[htbp]
\centering\small
\caption{Scenarios by task type and preregistered split. The split was stratified by family and pause pattern, not
by task type.}
\label{tab:tasksplit}
\begin{tabular*}{\textwidth}{@{\extracolsep{\fill}}l r r r@{}}
\toprule
Task type & training & test & total \\
\midrule
bugfix & 12 & 6 & 18 \\
feature & 10 & 3 & 13 \\
mixed & 18 & 10 & 28 \\
docs & 4 & 0 & 4 \\
explain & 1 & 2 & 3 \\
review & 2 & 2 & 4 \\
\bottomrule
\end{tabular*}

\end{table}

Long pauses and session length move the ratios less than the host does (\cref{tab:slicegap,tab:sliceturns}
in the appendix). On Opus, routed sessions with two long pauses came closer to break-even (shipped
\GapOpusShippedYes\X\ against \GapOpusShippedNo\X\ without), plausibly because the plain host also has to
re-write its expired cache after a pause. On Fable the ratios hardly move with pauses or length.

\subsection{Final hidden tests: a caution}\label{sec:finalpass}

\begin{table}[htbp]
\centering\small
\caption{Final hidden-test pass rate by arm, against the plain host on the same pairs, all splits. McNemar's
test compares the scenarios only the anchor passed with those only the arm passed.}
\label{tab:passrates}
\begin{tabular*}{\textwidth}{@{\extracolsep{\fill}}l l r r r r r@{}}
\toprule
Host & Arm & pairs & arm pass & anchor pass & anchor-only / arm-only & McNemar $p$ \\
\midrule
Fable 5.1 & sticky & 140 & 0.943 & 0.979 & 6 / 1 & 0.13 \\
Fable 5.1 & shipped & 140 & 0.936 & 0.979 & 7 / 1 & 0.07 \\
Fable 5.1 & sonnet & 140 & 0.943 & 0.979 & 5 / 0 & 0.06 \\
Fable 5.1 & second anchor (A/A) & 28 & 0.964 & 1.000 & 1 / 0 & 1.00 \\
\addlinespace[3pt]
Opus 5.5 & sticky & 140 & 0.936 & 0.957 & 5 / 2 & 0.45 \\
Opus 5.5 & shipped & 140 & 0.936 & 0.957 & 4 / 1 & 0.38 \\
Opus 5.5 & sonnet & 140 & 0.943 & 0.957 & 4 / 2 & 0.69 \\
Opus 5.5 & second anchor (A/A) & 28 & 0.929 & 1.000 & 2 / 0 & 0.50 \\
\bottomrule
\end{tabular*}

\end{table}

The turn-by-turn quality test is the preregistered one, and it passes. The final hidden tests tell a
slightly less comfortable story. On Fable, each routed arm lost more scenarios that the anchor passed than
it won: shipped \PassFableShippedAnchorOnly\ to \PassFableShippedArmOnly, sticky
\PassFableStickyAnchorOnly\ to \PassFableStickyArmOnly, Sonnet \PassFableSonnetAnchorOnly\ to
\PassFableSonnetArmOnly. \term{McNemar's test} asks whether such a lopsided split could be chance; here
$p$ ranges from \PassFableMinP\ to \PassFableMaxP, so it could. With so few discordant scenarios the test has
little power. Pooling the three routed Fable arms gives a clearer picture: they passed the final hidden tests
\FinalPoolAll\ points less often than the plain host across all scenarios (95\,\% interval
\FinalPoolAllCI\ points) and \FinalPoolTest\ points less often on the test scenarios (\FinalPoolTestCI).
This is exploratory and was not preregistered, but it is a real caution: a routed deployment should watch the
final pass rate.

\subsection{The noise floor}\label{sec:aa}

The \term{A/A} arm runs the plain host twice on the same scenario. Its ratio should be 1, and its spread
says how much two identical sessions differ by chance. The standard deviation of the A/A log ratio was about
\AAAllSD\ per pair, or \AAAllSDSession\ per session, far smaller than the effects above (\cref{tab:aa}). On
Fable the A/A ratio was \AAFable\ (\AAFableCI). On Opus it was \AAOpus\ (\AAOpusCI): the interval excludes 1. Two identical Opus
sessions differed by about \AAOpusPct\,\% on average, and we cannot explain why. A linear start-offset effect does not explain it: a
regression of the A/A log ratio on the start offset between the two copies gives a slope of
\RvAASlopeOpus\ per second on Opus (95\,\% CI \RvAASlopeOpusCI) and \RvAASlopeFable\ on Fable (\RvAASlopeFableCI),
and at zero offset the Opus ratio is still \RvAAInterceptOpus. The A/A log ratio does track the difference in
request counts between the two copies (correlation \RvAAReqCorrOpus\ on Opus, \RvAAReqCorrFable\ on Fable), which
suggests run-to-run variation in how the agent works, though a harness cause is not ruled out. With two A/A intervals, one
excluding 1 is not remarkable. If every Opus arm carried the same bias, removing it would move the Opus ratios from
\CfOpusStickyPairTwo--\CfOpusSonnetPairTwo\ to about \OpusAdjLow--\OpusAdjHigh, with the lowest lower bound at
\OpusAdjMinLower; the H2 verdicts would hold.

\subsection{Interim numbers versus final}\label{sec:interim}

\begin{figure}[htbp]
\centering
\pgfplotslegendfromname{interimlegend}\par\smallskip
% Cost ratios as the evidence accumulated: pilot, interim dashboard, final train, final test, confirmatory.
\begin{tikzpicture}
\begin{groupplot}[benchmark,
  group style={group size=2 by 1, horizontal sep=1.8cm},
  width=0.47\linewidth, height=6cm, xmin=-0.3, xmax=4.3,
  xtick={0,1,2,3,4}, xticklabels={Pilot, Interim, Final train, Final test, Confirm.},
  xticklabel style={rotate=35, anchor=north east, font=\small},
  ymajorgrids, title style={font=\small\bfseries},
  cycle list={{okblue, mark=*, solid}, {okorange, mark=square*, dashed}, {okgreen, mark=triangle*, densely dotted}},
  every axis plot/.append style={line width=1pt, mark size=2.4pt}]
\nextgroupplot[title={Fable host}, ylabel={Cost ratio vs plain host}, ymin=0.45, ymax=0.8,
  legend to name=interimlegend, legend columns=3, legend cell align=left,
  legend style={/tikz/every even column/.append style={column sep=10pt}}]
\addplot+ table[x=stage, y=Sticky]{generated/data/interim-fable.dat}; \addlegendentry{sticky (choose once)}
\addplot+ table[x=stage, y=Shipped]{generated/data/interim-fable.dat}; \addlegendentry{shipped (re-chooses each turn)}
\addplot+ table[x=stage, y=Sonnet]{generated/data/interim-fable.dat}; \addlegendentry{Sonnet only}
\nextgroupplot[title={Opus host}, ymin=0.95, ymax=1.75]
\draw[black!60, line width=0.6pt] (axis cs:-0.3,1) -- (axis cs:4.3,1);
\addplot+ table[x=stage, y=Sticky]{generated/data/interim-opus.dat};
\addplot+ table[x=stage, y=Shipped]{generated/data/interim-opus.dat};
\addplot+ table[x=stage, y=Sonnet]{generated/data/interim-opus.dat};
\end{groupplot}
\end{tikzpicture}

\caption{Cost ratio of each routed arm as the evidence accumulated: the screening pilot, the interim dashboard
during the run (partial training data), the final training and test splits (all cost-valid pairs) and the
confirmatory test result (pair reading).}
\label{fig:interim}
\howtoread{Read each line left to right. A flat line means the early numbers were already right. The Fable
lines are nearly flat. The Opus lines jump: the pilot overstated the extra cost and the interim dashboard
understated it.}
\end{figure}

On Fable, the pilot's ratios (sticky \PilotFableSticky, shipped \PilotFableShipped) were already close to the
final test values (\FinalTestFableSticky, \FinalTestFableShipped). On Opus, the pilot overstated the
penalty (sticky \PilotOpusSticky\ against \CfOpusStickyPairTwo\ in the confirmatory test) and the interim
dashboard understated it (\InterimOpusSticky). Small samples point in the right direction here but not at the
right size, which is why the decision waited for the preregistered test.

\subsection{Limits of the savings model}\label{sec:modellimits}

The model passed its preregistered check, and that check is about \emph{pooled} accuracy: are the intervals
right across all test pairs, and is the overall total inside its interval? It is not a check of each host and
arm separately, and the model's structure makes per-host figures unreliable. Each arm effect and each
task-type effect is shared by both hosts, although the data show the hosts behave very differently.
\begin{itemize}
  \item Across the \ModelCells\ host, task-type and arm cells, the model's point prediction lies outside the
    measured 90\,\% range in \ModelOutside\ (\cref{tab:per1000full}).
  \item It ranks the arms wrongly within a host. On the Fable test mix it predicts that Sonnet saves
    \$\TestMixFableSonnetModel\ per 1,000 sessions, less than shipped (\$\TestMixFableShippedModel); the
    measurement shows the opposite, \$\TestMixFableSonnetObs\ against \$\TestMixFableShippedObs.
\end{itemize}
Use the model for what it was validated for, pooled forecasts over a mix like this campaign's, and use the
measured tables for any per-host, per-task or per-arm figure.
